% auto-generated by make exhibits -- do not edit
\begin{table}[!htbp]\centering\footnotesize
\caption{The near-miss chain}\label{tab:chain}
\begin{tabular}{@{}lccc@{}}
\toprule
years since failure k & maj-fail issued $\le$k (n at risk) & min-fail issued $\le$k (n) & gap \\
\midrule
1 & 1.9\% (n=371) & 4.6\% (n=174) & -2.7\% \\
2 & 4.9\% (n=327) & 8.4\% (n=155) & -3.5\% \\
3 & 6.7\% (n=326) & 10.3\% (n=155) & -3.6\% \\
4 & 7.9\% (n=303) & 12.8\% (n=148) & -4.9\% \\
5 & 10.6\% (n=303) & 15.0\% (n=147) & -4.4\% \\
6 & 10.7\% (n=270) & 16.0\% (n=106) & -5.3\% \\
7 & 13.0\% (n=270) & 18.9\% (n=106) & -5.9\% \\
8 & 14.5\% (n=256) & 20.8\% (n=96) & -6.4\% \\
\bottomrule
\end{tabular}
\begin{minipage}{0.96\linewidth}\vspace{2pt}\scriptsize
\emph{Notes:} CA failed GO measures; per-cell observability restrictions; new-money issues only.
\end{minipage}
% BEGIN READING (strip for submission)
\begin{minipage}{0.96\linewidth}\vspace{4pt}\scriptsize
\textbf{Reading.} The near-miss deficit widens with the window: not a truncation artefact.
\end{minipage}
% END READING
\end{table}
